% !TEX encoding = UTF-8 Unicode
% !TEX root = thesis-main.tex

\chapterr{Εφαρμογή \textlatin{Android}}

Στο προηγούμενο κεφάλαιο παρουσιάσαμε και αναλύσαμε όλες τις λειτουργίες της πλατφόρμας του \textlatin{MeL} για κάθε επίπεδο χρήστη. Στο κεφάλαιο αυτό θα παρουσιάσουμε την ανάπτυξη της εφαρμογή του \textlatin{MeL} για κινητά τηλέφωνα με λειτουργικό σύστημα \textlatin{Android}. 


\section{Εισαγωγή στο \textlatin{Android}}

Το \textlatin{Android} αποτελεί ένα λειτουργικό σύστημα ανοιχτού κώδικα το οποίο χρησιμοποιείται τα τελευταία χρόνια σε ολοένα και αυξανόμενο αριθμό κινητών τηλεφώνων (\textlatin{smartphones}), καθώς και σε ηλεκτρονικές ταμπλέτες (\textlatin{tablet}). Όλες οι λειτουργίες του κινητού τηλεφώνου ελέγχονται από εφαρμογές (\textlatin{applications}) τις οποίες οι χρήστες μπορούν να κατεβάσουν και να εγκαταστήσουν ελεύθερα. Το γεγονός ότι αυτές οι συσκευές συνήθως τροφοδοτούνται από μια μπαταρία τις κάνει να απαιτούν από το λειτουργικό σύστημα που διαθέτουν να μην καταχράται τους πόρους του υλικού. Ο πυρήνας του \textlatin{Android} είναι βασισμένος στον πολύ γνωστό και αρκετά δοκιμασμένο πυρήνα του \textlatin{Linux}. Κληρονομώντας όλα τα χαρακτηριστικά ασφάλειας του \textlatin{Linux} και όλες τις διαχειριστικές τεχνικές μνήμης και επεξεργαστή που αυτό διαθέτει, το \textlatin{Android}, καθίσταται ένα αρκετά αξιόπιστο σύστημα.

\subsection{Αρχιτεκτονική}

Το \textlatin{Android} αποτελείται από μια στοίβα λογισμικών τμημάτων με ξεκάθαρους και αυστηρά καθορισμένους ρόλους όπως φαίνεται στο Σχήμα \ref{fig:android_archit} \cite{android}.

\begin{figure}[tbh]
	% \centerline{\includegraphics[scale=0.6]{figures/ch4-architecture_android}}
	\caption{Δομή της αρχιτεκτονικής \textlatin{Android}.}
	\label{fig:android_archit}
\end{figure}

Στο κατώτερο επίπεδο, βρίσκεται ο πυρήνας του \textlatin{Linux}. Ο πυρήνας του \textlatin{Linux} είναι γνωστός για την μεγάλη ποικιλία αρχιτεκτονικών επεξεργαστών με τις οποίες είναι συμβατός.

Στο επίπεδο ακριβώς πάνω από τον πυρήνα του \textlatin{Linux} βρίσκονται οι βασικές βιβλιοθήκες του συστήματος. Εδώ βρίσκονται οι υλοποιήσεις για βιβλιοθήκες γραφικών, για συστήματα Βάσεων Δεδομένων, η βασική βιβλιοθήκη της γλώσσας προγραμματισμού \textlatin{C} και άλλα. Όλες αυτές οι βιβλιοθήκες τρέχουν στον πυρήνα του \textlatin{Linux}. Το επίπεδο αυτό προσφέρει τα απαραίτητα στοιχεία ώστε να είναι δυνατή η υλοποίηση του \textlatin{Android Runtime}, του βασικού μηχανισμού, δηλαδή, που απαιτείται για την εκτέλεση των εφαρμογών που αναπτύσσονται για το περιβάλλον του \textlatin{Android}. 

Το \textlatin{Android Runtime} περιλαμβάνει την εικονική μηχανή (\textlatin{virtual machine}) που διερμηνεύει τις εφαρμογές \textlatin{Android} καθώς και την υλοποίηση των βασικών βιβλιοθηκών της \textlatin{Java} που προσφέρονται στον προγραμματιστή κατά την διαδικασία ανάπτυξης των δικών του εφαρμογών.

Βασισμένο στο \textlatin{Android Runtime} και στις βασικές βιβλιοθήκες του συστήματος, το προτελευταίο επίπεδο είναι το ονομαζόμενο \textlatin{Application Framework}. Το επίπεδο αυτό προσφέρει στον προγραμματιστή μια πληθώρα δυνατοτήτων σχετιζόμενων με το λειτουργικό σύστημα και την συσκευή που το φιλοξενεί στο σύνολο τους, όπως πρόσβαση σε διάφορα περιφερειακά της συσκευής, πρόσβαση σε λειτουργίες του ίδιου του λειτουργικού και άλλα.

Τέλος, το ανώτερο επίπεδο της στοίβας είναι το επίπεδο όπου βρίσκονται οι εφαρμογές για το περιβάλλον του \textlatin{Android}. Σε αυτό το επίπεδο αναπτύσσονται οι εφαρμογές και μπορούν να αξιοποιηθούν όλες οι δυνατότητες που προσφέρουν τα δυο προαναφερθέντα επίπεδα.\section{Εφαρμογή \textlatin{MeL}}

Στην προηγούμενη ενότητα παρουσιάσαμε κάποια εισαγωγικά στοιχεία που αφορούν το λειτουργικό σύστημα \textlatin{Android} και την αρχιτεκτονική του. Στην ενότητα αυτή θα παρουσιάσουμε τις λειτουργίες της εφαρμογής του \textlatin{MeL}. Η εφαρμογή αναπτύχθηκε με το πρόγραμμα \textlatin{Android Studio \cite{android_studio}} το οποίο κυκλοφόρησε το Δεκέμβριο του 2014 και από τότε αποτελεί το επίσημο ολοκληρωμένο περιβάλλον ανάπτυξης \textlatin{Android} εφαρμογών.

\subsection{Λειτουργίες εφαρμογής}


\section{Διαφορές εφαρμογής και πλατφόρμας \textlatin{MeL}}

Μετά την ανάλυση των λειτουργιών της εφαρμογής στην προηγούμενη ενότητα, στην ενότητα αυτή θα μιλήσουμε για τον σκοπό και τις διαφορές της με την πλατφόρμα του \textlatin{MeL}.

\begin{table}[t]
\centering
\begin{tabular}{|l||c|c|}
	\hline\rowcolor{gray} 
	ΛΕΙΤΟΥΡΓΙΕΣ & ΕΦΑΡΜΟΓΗ & ΠΛΑΤΦΟΡΜΑ \\
	\hline\hline
	Είσοδος/Αποσύνδεση &\checkmark & \checkmark \\
	Εγγραφή & - & \checkmark \\
	Αναζήτηση & - & \checkmark \\
	Επεξεργασία προφίλ & - & \checkmark \\
	Δημιουργία/Επεξεργασία Άρθρου/Τεστ & - & \checkmark \\
	Ανάγνωση Άρθρου & \checkmark & \checkmark \\
	Λύση Τεστ & \checkmark & \checkmark \\
	Κοινωνική Δικτύωση & - & \checkmark \\
	\hline
\end{tabular}
\caption{Διαφορές εφαρμογής και πλατφόρμας.}
\label{table:app_web}
\end{table}


\section{Σύνοψη}

Ο σκοπός της εφαρμογής του \textlatin{MeL} είναι να παρέχει εύκολη και γρήγορη πρόσβαση των χρηστών στην ηλεκτρονική μάθηση της πλατφόρμας. Στον Πίνακα \ref{table:app_web} παρουσιάζονται οι λειτουργίες της πλατφόρμας του συστήματος και ποιες από αυτές υποστηρίζονται από την εφαρμογή. Παρατηρούμε ότι οι λειτουργίες της εφαρμογής εστιάζουν κυρίως στο κομμάτι της ηλεκτρονικής μάθησης του \textlatin{MeL}, το οποίο αποτελεί και το βασικότερο σκοπό του συστήματος. Έτσι, οι λειτουργίες που μπορούν να εκτελέσουν οι χρήστες μέσα από την εφαρμογή σε σχέση με την πλατφόρμα, είναι η ανάγνωση άρθρου, η λύση τεστ ανάλογα με την κατηγορία που επιθυμούν και να ενημερωθούν για τα άρθρα ή τα τεστ που έχουν ανέβει πρόσφατα στην πλατφόρμα.







